%% ma: L12-B02-0010
%% lop: 12
%% chuong: 1
%% bai: 2
%% dang: 12.2.7
%% loai: DS
%% muc: [TH, TH, NB, TH]
%% dap_an: "ĐSĐS"
%% nguon: "Ảnh giáo viên gửi 11/10/2026 (Đề kiểm tra giữa HK1 Toán 12, THPT Việt Đức 2026-2027, mã đề 0103), Phần II Câu 3"
%% kien_thuc: "Bài 1 (cực trị), Bài 2 (giá trị nhỏ nhất trên một đoạn), phương trình đường thẳng qua hai điểm (lớp 10)"
%% trang_thai: chinh-thuc
%% ghi_chu: "Đã duyệt 12/10/2026."
\begin{ex}
Cho hàm số $f(x)=x^3-6x^2+9x+2$.
\choiceTF
{\True Đường thẳng đi qua hai điểm cực trị của đồ thị có hệ số góc bằng $-2$.}
{Đồ thị hàm số có điểm cực tiểu là điểm $I(1;6)$.}
{\True $f'(x)=3(x-1)(x-3)$.}
{Giá trị nhỏ nhất của hàm số $f(x)$ trên đoạn $[0;2]$ bằng $4$.}
\loigiai{$f'(x)=3x^2-12x+9=3(x-1)(x-3)$ nên c) đúng.

a) $f'(x)=0\Leftrightarrow x=1$ hoặc $x=3$; hai điểm cực trị là $(1;6)$ và $(3;2)$. Hệ số góc của đường thẳng qua hai điểm là $\dfrac{2-6}{3-1}=-2$. Đúng.

b) $f'(x)>0$ trên $(-\infty;1)$, $f'(x)<0$ trên $(1;3)$ nên $x=1$ là điểm cực đại, điểm cực đại của đồ thị là $(1;6)$; điểm cực tiểu là $(3;2)$. Sai.

d) $f(0)=2$, $f(1)=6$, $f(2)=4$ nên $\min_{[0;2]}f(x)=f(0)=2\ne4$. Sai.}
\end{ex}
